\documentclass[12pt]{article}
\usepackage[utf8]{inputenc}
\usepackage[T1]{fontenc}
\usepackage[margin=1in]{geometry}
\usepackage{setspace}
\usepackage{xcolor}
\usepackage[colorlinks=true, urlcolor=blue!60!black, linkcolor=blue!60!black]{hyperref}
\setlength{\parindent}{0em}
\setlength{\parskip}{0.6em}
\renewcommand{\baselinestretch}{1.0}
\begin{document}
\begin{center}
{\LARGE \textit{Griswold v. Connecticut}, 381 U.S. 479 (1965)\par}
\vspace{0.5em}
{\large Case Brief --- Two-Page Edition\par}
\end{center}

\section{Memory Jogger}
The Court struck down Connecticut's ban on contraceptive use by married couples, finding a right of marital privacy in the ``penumbras'' of the Bill of Rights (\href{https://supreme.justia.com/cases/federal/us/381/479/}{381 U.S. 479}).

\section{Detailed Case Facts}
Connecticut's 1879 statute made it a crime to use any drug or instrument to prevent conception, and an accessory statute reached anyone who assisted. After earlier challenges failed on standing and ripeness grounds, Estelle Griswold of Planned Parenthood and Dr. C. Lee Buxton opened a New Haven clinic in November 1961, advising married couples and prescribing contraceptives. They were arrested as accessories.

\section{Procedural History}
Griswold and Buxton were convicted and fined \$100 each; the Connecticut appellate courts affirmed. The Supreme Court noted probable jurisdiction and reversed.

\section{Judicial Votes}
Reversed, 7--2. Douglas wrote for Warren, Clark, Brennan, and Goldberg; Goldberg (with Warren and Brennan) concurred; Harlan and White concurred in the judgment; Black and Stewart dissented.

\section{Holding}
The accessories had standing to assert married couples' rights, and a ban on married couples' use of contraceptives violates a fundamental right of marital privacy protected by the penumbras of specific Bill of Rights guarantees.

\section{Analysis of Opinions}
Douglas disclaimed \textit{Lochner} and found ``zones of privacy'' in the First, Third, Fourth, Fifth, and Ninth Amendments; a ban on use, rather than sale, had a ``maximum destructive impact'' on marriage. Goldberg invoked the Ninth Amendment and the traditions of the people. Harlan grounded the right directly in due process and ordered liberty; White found the law irrational. Black insisted no specific provision barred the law, and Stewart called it ``an uncommonly silly law'' that only the legislature could repeal.

\section{Examples: Future Applications}
\textbf{Same-Side:} a ban on distributing contraceptives to unmarried adults fails, as in \textit{Eisenstadt}. \textbf{Opposite-Side:} a prescription requirement for hormonal contraceptives regulates distribution, not use, and is valid. \textbf{Fence-Sitter:} parental notification before prescribing birth control to minors under 16.

\section{Critique}
Robert Bork attacked the leap from specific guarantees to a freestanding right of privacy, and Justice Thomas in \textit{Dobbs} urged reconsidering \textit{Griswold}. Feminist scholars note that the marital-privacy frame excluded unmarried women, and later Courts replaced penumbras with Harlan's due process approach.

\section{Key Quotations}
\begin{quote}
``The foregoing cases suggest that specific guarantees in the Bill of Rights have penumbras, formed by emanations from those guarantees that help give them life and substance.''
\end{quote}
\noindent --- Justice Douglas, Opinion of the Court, 381 U.S. at 484.

\end{document}
